Na presente seção será apresentada uma contextualização do tema, posteriormente a justificativa e objetivos. Por fim, a estrutura do trabalho é mostrada.

\section{Contextualização do tema}
A agricultura é uma atividade fundamental para a economia global, sendo responsável pela produção de alimentos e matérias-primas essenciais para a sobrevivência e o bem-estar da humanidade. De acordo \textcite{carmody_fao2023}, o setor agrícola enfrenta desafios significativos. A crescente demanda por alimentos devido ao aumento populacional e as mudanças climáticas afetam a disponibilidade de recursos naturais. Diante desses desafios, a gestão eficiente dos recursos hídricos e a adoção de tecnologias avançadas tornam-se cruciais para garantir a sustentabilidade e a produtividade das culturas \parencite{carmody_fao2023}.

No Brasil, a agricultura desempenha um papel vital na economia, sendo um dos principais motores de crescimento e desenvolvimento do país \parencite{Ramos_irrigacao2022}. O Brasil é um dos maiores produtores e exportadores de alimentos do mundo, destacando-se na produção de soja, milho, café e carne bovina. 

\textcite{Ramos_irrigacao2022} apresenta a irrigação como uma prática amplamente estudada no país para aumentar a produtividade agrícola e garantir a segurança alimentar. Métodos tradicionais de irrigação muitas vezes carecem de precisão e eficiência, levando ao desperdício de água e outros recursos \parencite{Pereira_irrigation2002}.

A evapotranspiração (ET) é uma das principais variáveis a serem consideradas na gestão eficiente da água \parencite{yang_nature2023}. Pois, a ET é fundamental para determinar as necessidades hídricas das culturas e otimizar o uso da água na agricultura \parencite{carmody_fao2023, yang_nature2023}. Esta variável (ET) representa a perda de água do solo por evaporação e pela transpiração das plantas. 

O desenvolvimento de sistemas de monitoramento em tempo real para a irrigação surge como uma abordagem promissora para otimizar o uso da água e melhorar o rendimento das culturas \parencite{Vijh_system2024, Dutta_system2024, Amine_system2024}. De acordo com \textcite{Prathibha_system2017}, tais sistemas fornecem aos agricultores informações precisas e atualizadas sobre as condições do solo e das plantas, permitindo uma gestão mais eficiente da irrigação ao longo do tempo.


\section{Justificativa}
Este estudo visa contribuir para o avanço do conhecimento na área da agricultura inteligente, fornecendo uma solução prática e eficaz para o monitoramento em tempo real da irrigação. A implementação de tecnologias de monitoramento e controle precisos pode melhorar significativamente a sustentabilidade e a eficiência dos sistemas de produção agrícola.

Ao fornecer aos agricultores ferramentas que lhes permitam monitorar e controlar o processo de irrigação de forma mais precisa e eficiente, espera-se contribuir para a redução do desperdício de água. Isso, por sua vez, pode resultar na otimização do uso de recursos e no aumento da produtividade agrícola.

A relevância deste estudo se evidencia pela necessidade de soluções inovadoras que abordem os desafios enfrentados pela agricultura moderna \parencite{carmody_fao2023}. Em especial, a irrigação, cujo consumo global e nacional de água é significativo \parencite{Ramos_irrigacao2022}. 

\section{Objetivos}

Este trabalho tem como objetivo geral desenvolver um sistema de monitoramento em tempo real para o cultivo agrícola. A proposta é utilizar tecnologias modernas e o método científico da ETc, visando aumentar a eficácia do monitoramento da cultura e de suas necessidades hídricas.

Os objetivos específicos deste trabalho são os seguintes:

\begin{itemize}
    \item Implementar um sistema de telemetria de sensores para coletar dados ambientais-chave em tempo real;
    \item Integrar algoritmos de cálculo da ETc ao sistema para fornecer métricas;
    \item Desenvolver uma interface de usuário intuitiva para facilitar o acesso e a compreensão dos dados pelos agricultores.
\end{itemize}

\section{Estrutura do trabalho}
Este trabalho está estruturado da seguinte forma:

\begin{itemize}
    \item O Capítulo 1 apresenta esta introdução, contextualizando o tema, justificando sua relevância, delineando os objetivos e descrevendo a metodologia adotada;
    \item O Capítulo 2 revisa a literatura relacionada, fornecendo uma base teórica para o desenvolvimento do sistema proposto;
    \item O Capítulo 3 descreve em detalhes a metodologia de desenvolvimento do sistema;
    \item O Capítulo 4 apresenta os resultados dos testes e validações realizados em campo;
    \item Por fim, o Capítulo 5 conclui o trabalho, discutindo suas contribuições, limitações e possíveis direções futuras.
\end{itemize}